\par\bigskip\pdfbookmark[0]{Über den Zusammenhang. Fragment IV}{nexus-fragment-IV}\centerline{[IV.]}\medskip
Man bezeichne mit $(n)$ die Anzahl der Werte $x$ aus der Gesamtheit
\[
0,1,2,3,4\dotsfour p^n-1
\]
für die $xx-D=xx-ap^\mu$ durch $p^n$ teilbar ist.
\[
\begin{array}{c@{\qquad}cc}
1)\ \mu\text{ ungerade z. B. }=7.&\multicolumn2c{2)\ \mu\text{ gerade z. B. }=6}\\
(1)=1&aNp&aRp\\
(2)=p&(1)=1&(1)=1\\
(3)=p&(2)=p&(2)=p\\
(4)=pp&(3)=p&(3)=p\\
(5)=pp&(4)=pp&(4)=pp
\end{array}
\]

\SourceLeaf{294}
\[
\begin{array}{c@{\qquad}c@{\qquad}c}
(6)=p^3&(5)=pp&(5)=pp\\
(7)=p^3&(6)=p^3&(6)=p^3\\
(8)=0&(7)=0&(7)=2p^3\\
(9)=0&(8)=0&(8)=2p^3\\
\text{usw.}&\text{usw.}&\text{usw.}
\end{array}
\]

\[
\begin{array}{c@{\qquad}c}
\text{Man mache nun}&\text{Dann ist}\\[.4em]
(1)-\dfrac{(2)}p=(1)'&fp=(1)'\\[.4em]
(2)-\dfrac{(3)}p=(2)'&fpp=1+(2)'\\[.4em]
(3)-\dfrac{(4)}p=(3)'&fp^3=(1)'+(3)'\\[.4em]
(4)-\dfrac{(5)}p=(4)'&fp^4=1+(2)'+(4)'\\
\text{usw.}&\text{usw.}
\end{array}
\]
Setzt man also $\frac{p-1}p(1+\frac{fp}p+\frac{fpp}{pp}+\frac{fp^3}{p^3}+\text{usw.})=T$, so ist
\[
\begin{aligned}
\frac{p+1}p T&=1+\frac{(1)'}p+\frac{(2)'}{pp}+\frac{(3)'}{p^3}+\frac{(4)'}{p^4}+\text{usw.}\\
&=1+\frac{(1)}p=1+\frac1p
\end{aligned}
\]
Also $T=1$
\BookRule{.30\textwidth}
